\chapter{Tips and examples}
\label{ch:tips-and-examples}

This chapter presents some basic tips and a few examples on how to use \LaTeX.

\section{Basic tips}
\label{sec:basic-tips}

\begin{itemize}
\item
The \verb+\\+ command has the same effect as \verb+\newline+.
\item
The \verb+\cleardoublepage+ command forces the next content to start in an odd page.
\item
The tilde character (\verb+~+) inserts a non-breaking space. Use it before citing a reference to avoid breaking the line: \verb+an example~\cite{label}+.
\item
The current font size is \myfontsize.
\item
Use the \verb+longtable+ environment for tables spanning multiple pages.
\item
The grave accent \`{} and the apostrophe \textquotesingle\ are the correct symbols to make quotations: ``this is an example''.
\end{itemize}

\section{Font styles}
\label{sec:font-styles}

\begin{itemize}
\item
\verb+\textnormal{}+ --- \textnormal{document font family}.
\item
\verb+\textrm{}+ --- \textrm{roman font family}.
\item
\verb+\textsf{}+ --- \textsf{sans serif font family }.
\item
\verb+\texttt{}+ --- \texttt{teletypefont family}.
\item
\verb+\textit{}+ --- \textit{italic shape}.
\item
\verb+\textsl{}+ --- \textsl{slanted shape}.
\item
\verb+\textsc{}+ --- \textsc{small capitals}.
\item
\verb+\textbf{}+ --- \textbf{bold}.
\item
\verb+\mybold{}+ --- \mybold{another type of bold}. (The \verb+\mybold+ command is defined in the \verb+tex/config/after-begin.tex+ file.)
\end{itemize}

\section{Colors}
\label{sec:colors}

{\color{red} Red colored text} from the \texttt{color} package.
{\color{NavyBlue} And NavyBlue colored text} from the \texttt{xcolor} package.

\section{Footnotes}
\label{sec:footnotes}

This is a labeled footnote\cref{foot:example}. A footnote can be referenced multiple times\footnote{\label{foot:example}This is a footnote example.}. Again, the same footnote is referenced\cref{foot:example}.

\subsection{A table with a footnote}
\label{ssec:a-table-with-a-footnote}

A table footnote is used in \Cref{tab:table-with-footnote}.

\FloatBarrier
\begin{table}[!htbp]
\begin{center}
\label{tab:table-with-footnote}
\caption{A table with a \texttt{\textbackslash tablefootnote}.}
\vspace*{5pt}
\begin{tabular}{||m{3cm}|A{3cm}|S{3cm}|D{3cm}||N}
\hline
justified \dummytext & left-aligned \dummytext & centered \dummytext & right-aligned \dummytext & \\
\hline
A footnote\tablefootnote{This is a table footnote.} & 2 & 3 & 4 & \\
\hline
Single cell & \multicolumn{3}{c||}{A multi-column cell} & \\
\hline
\multirow{2}[1]{*}{A multi-row cell} & A simple & ... & ... & \\
& example & ... & ... & \\
\hline
\end{tabular}
\end{center}
\end{table}
\FloatBarrier


\section{Abbreviations}
\label{sec:abbreviations}

\verb+\gls{label}+ and \verb+\glslink{label}{text}+ are two possible commands for making use of abbreviations. For example, the commands \verb+\gls{afk}+ (first call), \verb+\gls{afk}+ (second call), and \verb+\glslink{afk}{insert specific text}+ produce respectively ``\gls{afk}'', ``\gls{afk}'' and ``\glslink{afk}{insert specific text}''.

A list of some commands follow.

\begin{itemize}
\item
\verb+\gls{afk}+ produces ``\gls{afk}''.
\item
\verb+\glslink{afk}{text}+ produces ``\glslink{afk}{text}''.
\item
\verb+\acrshort{afk}+ produces ``\acrshort{afk}''.
\item
\verb+\acrlong{afk}+ produces ``\acrlong{afk}''.
\end{itemize}

Other abbreviations include: good work (\acrshort{gw}); have fun (\acrshort{hf}); good work and have fun (\acrshort{gwhf}). Note that the latter contains nested abbreviations.

\section{Equations}
\label{sec:equations}

\Cref{eq:example1} is a numbered equation.

\begin{equation}
\label{eq:example1}
x = 1 + y
\end{equation}

The following equation is not numbered, and thus cannot be referenced.

\begin{equation*}
y = \sum_{i=1}^{N}{x_i}
\end{equation*}

The \verb+\myscore+ pre-defined math function is used by the \Cref{eq:example2}.

\begin{equation}
\label{eq:example2}
\myscore(d) = \frac{1}{d^2}
\end{equation}

\section{Figures}
\label{sec:figures}

An example of a figure is shown in \Cref{fig:aveiro}. The \verb+\fbox+ command \fbox{draws a box} around its content.

\begin{figure}[!htbp]
\centering
\fbox{\includegraphics[width=0.1\textwidth]{img/logos/aveiro.pdf}}
\caption[This caption is shown at the index.]{This caption is shown below the figure.}
\label{fig:aveiro}
\end{figure}


\section{Rotating pages}
\label{sec:rotating-pages}

Horizontal pages can be obtained using the \verb+landscape+ environment  from the \verb+pdflscape+ package:

\begin{verbatim}
\begin{landscape}
Add some text or figures.
\end{landscape}
\end{verbatim}

\section{A long table}
\label{sec:a-long-table}

\Cref{tab:a-long-table} presents a long table using the \verb+longtable+ package. This table can span multiple pages. The \verb+\afterpage+ commands forces the table to start at the top of a page.

\afterpage{\begingroup

\normalfont
\color{black}

% O   + p    + I*2    + p
% 0.0 + 0.36 + 0.02*2 + 0.60 = 1.0
\begin{longtable}[c]{Op{0.36\textwidth}IIp{0.60\textwidth}N}

\caption{A long table.}
\label{tab:a-long-table}\\\hline\hline

\textbf{Column 1} & \textbf{Column 2} &\\\hline\hline

0 & \dummytext\dummytext\dummytext\dummytext\dummytext & \\\hline
1 & \dummytext\dummytext\dummytext\dummytext\dummytext & \\\hline
2 & \dummytext\dummytext\dummytext\dummytext\dummytext & \\\hline
3 & \dummytext\dummytext\dummytext\dummytext\dummytext & \\\hline
4 & \dummytext\dummytext\dummytext\dummytext\dummytext & \\\hline
5 & \dummytext\dummytext\dummytext\dummytext\dummytext & \\\hline
6 & \dummytext\dummytext\dummytext\dummytext\dummytext & \\\hline
7 & \dummytext\dummytext\dummytext\dummytext\dummytext & \\\hline
8 & \dummytext\dummytext\dummytext\dummytext\dummytext & \\\hline
9 & \dummytext\dummytext\dummytext\dummytext\dummytext & \\\hline

\end{longtable}

\endgroup
}
